% Generated draft: requires mathematical and source review.
\documentclass[11pt]{article}
\usepackage[margin=25mm]{geometry}
\usepackage{fontspec}
\setmainfont{Latin Modern Roman}
\usepackage{amsmath,amssymb,mathtools,unicode-math}
\setmathfont{Latin Modern Math}
\usepackage{enumitem,needspace,xurl,hyperref}
\setlength{\parindent}{0pt}
\setlength{\parskip}{5pt}
\setlength{\emergencystretch}{3em}
\newcommand{\N}{\mathbb{N}}
\newcommand{\Z}{\mathbb{Z}}
\newcommand{\Q}{\mathbb{Q}}
\newcommand{\R}{\mathbb{R}}
\newcommand{\C}{\mathbb{C}}
\newcommand{\F}{\mathbb{F}}
\newcommand{\A}{\mathbb{A}}
\newcommand{\PP}{\mathbb{P}}
\newcommand{\E}{\mathbb{E}}
\newcommand{\eps}{\varepsilon}
\newcommand{\dd}{\,\mathrm{d}}
\begin{document}
\section*{UnsolvedMath ID: 801086}
\textbf{Source problem code:} catalogue entry
\subsection{Definitions and mathematical statement}
Let \(d\) be a positive integer. A polynomial automorphism of \(\C^d\) is a polynomial map \(f=(f_1,\ldots,f_d):\C^d\to\C^d\) admitting a polynomial inverse. Write \(\deg(f)=\max_i\deg(f_i)\), using total polynomial degree. With the usual first dynamical degree convention, define
\[
\lambda_1(f)=\lim_{n\to\infty}\bigl(\deg(f^{\circ n})\bigr)^{1/n},
\]
where \(f^{\circ n}\) denotes the \(n\)-fold iterate. A weak Perron number is a real algebraic integer \(\lambda\geq1\) such that every Galois conjugate \(\mu\) of \(\lambda\) satisfies \(|\mu|\leq\lambda\). Its algebraic degree is \([\Q(\lambda):\Q]\).

The supplied Question 38 asks whether, for every polynomial automorphism \(f\) of \(\C^d\), the number \(\lambda_1(f)\) is a weak Perron number with
\[
[\Q(\lambda_1(f)):\Q]\leq d-1.
\]
Here the displayed field degree is asserted only as part of the proposed conclusion that \(\lambda_1(f)\) is algebraic. The source does not specify the dimension range. In particular, the bound becomes zero for \(d=1\); whether the intended question assumes \(d\geq2\) requires clarification. The source also does not explicitly define its dynamical degree; the formulation above uses the first dynamical degree.

The supplied background reports computations by Blanc and Van Santen concerning cubic automorphisms of \(\C^3\) and automorphisms obtained by composing an affine transformation and a triangular transformation. These are motivation for the question, not restrictions on its stated universal quantifier.
\subsection{Short English statement}
The question asks whether the degree growth rate of every polynomial automorphism of \(\C^d\) is an algebraic integer dominating the absolute values of all its conjugates, with algebraic degree at most \(d-1\). The intended dimension range requires clarification. The supplied registry labels the problem open; this draft does not independently verify that status.
\subsection{Sources}
\begin{itemize}
\item Local UnsolvedMath registry, permanent ID \(801086\), source code \texttt{catalogue entry}, title \emph{Henon Maps: Open Problems --- Question 38}. The supplied record attributes the source list to Berger and Bedford (2023) and labels the problem open.
\item Berger and Bedford, \emph{Henon Maps: A List of Open Problems} (2023), Question 38: \url{https://arxiv.org/abs/2312.03907}. The supplied provenance records access on 2026-10-16 and extraction by \texttt{arxiv\_latex}. These are supplied provenance details, not a claim of independent source access or verification.
\item The supplied statement attributes the motivating computations and resulting question to Blanc and Van Santen, citing the source labels \texttt{[blanc-V1,blanc-V2]}. Full bibliographic entries for those labels were not supplied.
\end{itemize}
\end{document}
